\scanpage{29}

\uline{Cas particulier.} Soit $\deleted{\unsure{\mathcal F|F}}\,\deleted{\unsure{\mathcal F}}F$ un faisceau en groupes (pas néc. \unsure{comm|comm.}) sur $X$. Soit $\mathcal F$ la catégorie fibrée (sur les ouverts de $X$) des fibrés principaux homogènes sous $F$. Soit $Y$ une partie fermée de $X$. Alors

% slanted
\margin{Pas d'hyp. \unsure{noethérienne|noetherienne} sur $X$}

\uline{Proposition} 9\quad \deleted{Pour que} \inserted{Conditions équivalentes}

(i) $\mathcal F$ satisfait $P_1$ ($P_2$) rel. à $Y$

(ii) $F$ satisfait $P_1$ ($P_2$) rel. à $Y$.

(iii) On a $\underline H^0_Y(F)=0$ \deleted{pour $i<1$ ($i<2$) i.e. $i=0$ ($i=0,1$)} \inserted{($\underline H^0_Y(F)=0$ et $\underline H^1_Y(F)=0$)}

En effet, \deleted{on aura} si $\xi_0$ est le fibré trivial sur $X$, alors $F\simeq\underline{\mathrm{Hom}}(\xi_0,\xi_0)$, d'où $(\mathrm{i})\Rightarrow(\mathrm{ii})$. D'autre part, si $\xi,\eta\in\mathrm{Ob}\,\mathcal F(V)$, alors $\underline{\mathrm{Hom}}(\xi,\eta)$ est \uline{loc isom} (comme faisceau d'ens) à $F$, d'où $(\mathrm{ii})\Rightarrow(\mathrm{i})$.

Enfin $(\mathrm{ii})\Leftrightarrow(\mathrm{iii})$ par définition de $\underline H^0_Y(F)$ et $\underline H^1_Y(F)$ comme Ker et Coker de $F\unsure{\longrightarrow|-}i_*i^*F$.

% archive sheet number: 28
